The \code{BPatch\_process} class represents a running process, which
includes one or more threads of execution and an address
space.

\begin{apient}
  bool stopExecution()
  bool continueExecution()
  bool terminateExecution
\end{apient}


\apidesc{
  These three functions change the running state of the process.
  stopExecution puts the process into a stopped state.

  continueExecution continues execution of the process.
  terminateExecution terminates execution of the process and will
  invoke the exit callback if one is registered. Each function
  returns true on success, or false for failure. Stopping
  or continuing a terminated thread will fail and these functions
  will return false.
}

\begin{apient}
  bool isStopped()
  int stopSignal()
  bool isTerminated()
\end{apient}

\apidesc{
  These three functions query the status of a process. isStopped
  returns true if the process is currently stopped. If the
  process is stopped (as indicated by isStopped), then stopSignal can
  be called to find out what signal caused the
  process to stop. isTerminated returns true if the process has
  exited. Any of these functions may be called multiple
  times, and calling them will not affect the state of the process.
}

\begin{apient}
  BPatch_variableExpr *getInheritedVariable(BPatch_variableExpr &parentVar)
\end{apient}

\apidesc{
  Retrieve a new handle to an existing variable (such as one created
  by \code{BPatch\_process\:\:malloc}) that was created in a
  parent process and now exists in a forked child process. When a
  process forks all existing \code{BPatch\_variableExpr}s are
  copied to the child process, but the Dyninst handles for these
  objects are not valid in the child \code{BPatch\_process}.
  This function is invoked on the \code{BPatch\_process} of the child process,
  parentVar is a variable from the
  parent process, and a handle to a variable in the child process is
  returned. If parentVar was not allocated in the
  parent process, then NULL is returned.
}


\begin{apient}
  BPatchSnippetHandle *getInheritedSnippet(BPatchSnippetHandle &parentSnippet)
\end{apient}

\apidesc{
  This function is similar to getInheritedVariable, but operates on
  BPatchSnippetHandles. Given a child process that was
  created via fork and a BPatch­SnippetHandle, parentSnippet, from
  the parent process, this function will return a handle
  to parentSnippet that is valid in the child process. If it is
  determined that parentSnippet is not associated with
  the parent process, then NULL is returned.
}

\begin{apient}
  void detach(bool cont)
\end{apient}

\apidesc{
  Detach from the process. The process must be stopped to call this
  function. Instrumentation and other changes to the
  process will remain active in the detached copy. The cont parameter
  is used to indicate if the process should be
  continued as a result of detaching.
  Linux does not support detaching from a process while leaving it
  stopped. All processes are continued after detach on
  Linux.
}

\begin{apient}
  int getPid()
\end{apient}

\apidesc{
  Return the system id for the mutatee process. On UNIX based systems
  this is a PID. On Windows this is the HANDLE
  object for a process.
}

\begin{apient}
  enum BPatch_exitType {
    NoExit,
    ExitedNormally,
    ExitedViaSignal
  }
\end{apient}

\apidesc{
  The mechanism by which the process exited.
}

\begin{apient}
  BPatch_exitType terminationStatus()
\end{apient}

\apidesc{
  If the process has exited, terminationStatus will indicate whether
  the process exited normally or because of a signal.

  if the process was not a child of the Dyninst mutator; in this
  case, ExitedNormally will be returned in both normal and
  signal exit cases.
}

\begin{apient}
  int getExitCode()
\end{apient}

\apidesc{
  If the process exited in a normal way, getExitCode will return the
  associated exit code. Prior to Dyninst 8.2,
  getExitCode would return the argument passed to exit or the value
  returned by main; in Dyninst 8.2 and later, it
  returns the actual exit code as provided by the debug interface and
  seen by the parent process. In particular, on
  Linux, this means that exit codes are normalized to the range 0-255.
}

\begin{apient}
  int getExitSignal()
\end{apient}

\apidesc{
  If the process exited because of a received signal, getExitSignal
  will return the associated signal number.
}

\begin{apient}
  void oneTimeCode(const BPatch_snippet &expr)
\end{apient}

\apidesc{
  Cause the snippet expr to be executed by the mutatee immediately.
  If the process is multithreaded, the snippet is run on
  a thread chosen by Dyninst. If the user requires the snippet to be
  run on a particular thread, use the \code{BPatch\_thread}
  version of this function instead. The process must be stopped to
  call this function. The behavior is synchronous;
  oneTimeCode will not return until after the snippet has been run in
  the application.
}

\begin{apient}
  bool oneTimeCodeAsync(const BPatch_snippet &expr, void *userData = NULL)
\end{apient}

\apidesc{
  This function sets up a snippet to be evaluated by the process at
  the next available opportunity. When the snippet
  finishes running Dyninst will callback any function registered
  through \code{BPatch::registerOneTimeCodeCallback}, with
  userData passed as a parameter. This function return true on
  success and false if it could not post the oneTimeCode.
  If the process is multithreaded, the snippet is run on a thread
  chosen by Dyninst. If the user requires the snippet to
  be run on a particular thread, use the \code{BPatch\_thread} version of
  this function instead. The behavior is asynchronous;
  oneTimeCodeAsync returns before the snippet is executed.
  If the process is running when oneTimeCodeAsync is called, expr
  will be run immediately. If the process is stopped,
  then expr will be run when the process is continued.
}

\begin{apient}
  void getThreads(std::vector<BPatch_thread *> &thrds)
\end{apient}

\apidesc{
  Get the list of threads in the process.
}

\begin{apient}
  bool isMultithreaded()
  bool isMultithreadCapable()
\end{apient}

\apidesc{
  The former returns true if the process contains multiple threads;
  the latter returns true if the process can create
  threads (e.g., it contains a threading library) even if it has not yet.
}
